\documentclass[12pt]{article}
\usepackage[utf8]{inputenc}
\usepackage[T1]{fontenc}
\usepackage[margin=1in]{geometry}
\usepackage{setspace}
\usepackage{xcolor}
\usepackage[colorlinks=true, urlcolor=blue!60!black, linkcolor=blue!60!black]{hyperref}
\setlength{\parindent}{2em}
\setlength{\parskip}{1.25em}
\renewcommand{\baselinestretch}{1.0}
\begin{document}
\begin{center}
{\LARGE \textit{DHS v. Thuraissigiam}, 591 U.S. 103 (2020)\par}
\vspace{0.8em}
{\large Case Brief --- Four-Page Edition\par}
\end{center}
\vspace{1em}

\section{Memory Jogger}

A Sri Lankan asylum seeker caught 25 yards inside the border could not use habeas to challenge his negative credible-fear finding: the Court ruled against him 7--2. The five-Justice majority held that the Suspension Clause did not require habeas review to give him another chance to obtain admission and that a noncitizen at the threshold of entry has only the admission procedures Congress grants (\href{https://supreme.justia.com/cases/federal/us/591/19-161/}{591 U.S. 103}).

\section{Detailed Case Facts}

Under IIRIRA (1996), ``applicants for admission'' without valid documents who have not been present for two years face expedited removal unless they show a ``credible fear of persecution''---a ``significant possibility'' of asylum eligibility---before an asylum officer, with supervisory and \textit{de novo} immigration-judge review. Section 1252(e)(2) limits habeas review to whether the petitioner is an alien, was ordered removed, or already holds lawful permanent resident, refugee, or asylee status; courts may not review the merits of a negative credible-fear finding.

In January 2017 Vijayakumar Thuraissigiam, a Sri Lankan national, crossed the border at night without inspection and was caught 25 yards inside it. In his interview he said unidentified men had abducted and beaten him, but he did not link the attack to a protected ground. An asylum officer, a supervisor, and an immigration judge all found no credible fear.

\section{Procedural History}

Thuraissigiam petitioned for habeas in the Southern District of California, alleging for the first time persecution for Tamil ethnicity and political opinion and seeking a new chance to apply for asylum, not release from custody. The district court dismissed under \S\ 1252(e)(2). The Ninth Circuit reversed, holding the statute violated the Suspension and Due Process Clauses as applied, splitting with the Third Circuit. The Supreme Court reversed on June 25, 2020.

\section{Judicial Votes}

Reversed and remanded with instructions to dismiss, \textbf{7--2}.
\begin{itemize}
\item \textbf{Majority:} Justice Alito, joined by Chief Justice Roberts and Justices Thomas, Gorsuch, and Kavanaugh.
\item \textbf{Concurrences:} Justice Thomas; Justice Breyer, joined by Justice Ginsburg (in the judgment).
\item \textbf{Dissent:} Justice Sotomayor, joined by Justice Kagan.
\end{itemize}

\section{Holding}

Section 1252(e)(2) violates neither the Suspension Clause nor the Due Process Clause as applied. The writ as understood in 1789 secured release from unlawful executive detention, not administrative review or authorization to enter and remain. A noncitizen caught 25 yards inside the border stands ``on the threshold of initial entry'' and has only the procedural rights regarding admission that Congress confers.

\section{Analysis of Opinions}

\textbf{Alito (majority).} Taking ``the writ as it existed in 1789'' as the floor (\textit{INS v.\ St.\ Cyr}), Alito relied on Blackstone and Story: habeas tests custody and secures simple release, which the government was willing to provide by flying him to Sri Lanka. He sought something like an injunction or mandamus, as in \textit{Munaf v.\ Geren}. \textit{Somerset v.\ Stewart} and the deserting-sailor cases ordered release; any ability to stay was a collateral consequence. The ``finality era'' cases from \textit{Nishimura Ekiu} onward reviewed legal questions under the 1867 habeas statute, not the Suspension Clause. \textit{Boumediene} concerned release from military custody. On due process, under \textit{Knauff} and \textit{Mezei}, whatever procedure Congress authorizes is due process for one seeking initial entry; a person caught just after crossing has not ``effected an entry,'' and giving clandestine entrants more rights than port-of-entry applicants would create a perverse incentive.

\textbf{Thomas (concurring).} ``Suspension'' originally meant authorizing detention without bail or trial on mere suspicion; IIRIRA requires an inadmissibility finding and allows some review.

\textbf{Breyer (concurring in the judgment).} The majority should not have reached due process. Habeas is an ``adaptable remedy,'' and the broad rule could endanger long-term residents or citizens, but Thuraissigiam had never lived here and raised factual and procedural quibbles finality-era courts never reviewed.

\textbf{Sotomayor (dissenting).} The Court ``handcuff[ed]'' judicial oversight and recast legal challenges as a plea for admission. \textit{St.\ Cyr}, \textit{Boumediene}, \textit{Somerset}, and \textit{Heikkila v.\ Barber} show habeas tests the legality of removal orders, and the entry fiction now reaches people on U.S. soil with no limiting principle.

\textbf{Later developments.} \textit{Jones v.\ Hendrix} (2023) extended the 1789-focused method to bar successive \S\ 2255 claims.

\section{Examples: Future Applications}

\begin{enumerate}
\item \textbf{Same-Side (No Review): Rio Grande Crossing.} A migrant caught 100 yards from the river two hours after crossing challenges the officer's ``particular social group'' analysis and seeks a new interview. She seeks admission, not release, so the court must dismiss.
\item \textbf{Same-Side (No Review): Port of Entry.} An applicant who presented herself at a port of entry complains of an interpreter's dialect. She has never entered and has only the process Congress provides.
\item \textbf{Opposite-Side (Review): Interior Detention.} DHS holds an eight-year Houston resident, or a citizen, indefinitely without bond under a statute stripping all habeas review. He seeks simple release, the 1789 core of the writ.
\item \textbf{Opposite-Side (Review): Citizenship Claim.} A man in expedited removal offers a Texas birth certificate. Section 1252(e)(2) itself permits review of alienage.
\item \textbf{Fence-Sitter: Interior Expedited Removal.} DHS applies expedited removal to a Chicago resident of eighteen months. The Suspension Clause holding cuts against him, but the Court never said when presence matures into an entry (see \textit{Zadvydas}).
\end{enumerate}

\section{Critique}

\textbf{Anachronism and line-drawing.} Demanding a 1789 precedent for review of asylum screenings asks the impossible, since no immigration restrictions or asylum statutes existed then, and it sits uneasily with \textit{Boumediene}'s functionalism. Holding that 25 yards inside the border is not an entry, without defining when presence matures into constitutional protection, invites litigation over interior expedited removal.

\textbf{Unchecked error.} Insulating credible-fear determinations from review even for legal error risks \textit{refoulement} of bona fide refugees in high-volume, summary proceedings.

\section{Key Quotations}

\begin{quote}
``For these reasons, an alien in respondent's position has only those rights regarding admission that Congress has provided by statute.''
\end{quote}
\noindent --- Justice Alito, Opinion of the Court, 591 U.S. at 140.

\begin{quote}
``Habeas has traditionally been a means to secure release from unlawful detention, but respondent invokes the writ to achieve an entirely different end, namely, to obtain additional administrative review of his asylum claim and ultimately to obtain authorization to stay in this country.''
\end{quote}
\noindent --- Justice Alito, Opinion of the Court, 591 U.S. at 107.

\begin{quote}
``Today's decision handcuffs the Judiciary's ability to perform its constitutional duty to safeguard individual liberty and dismantles a critical component of the separation of powers. It will leave significant exercises of executive discretion unchecked in the very circumstance where the writ's protections `have been strongest.'\,''
\end{quote}
\noindent --- Justice Sotomayor, dissenting, 591 U.S. at 159.

\end{document}
